% Ad Atticum 2.1 (ad-atticum-02-01) --- English translation
% Translated by Claude (Anthropic) from the Latin (Perseus).
% Style guide: STYLE.md.

\ciceroLetterOpener{Cicero to Atticus, greetings}

\ciceroSection{1} On the Kalends of June, as I was setting out for Antium and gladly leaving the gladiatorial games of Marcus Metellus behind, your boy met me. He delivered me a letter from you, and the memoir of my consulship written in Greek. Reading it I was glad that I had given to Lucius Cossinius, some good while before, a book of mine on the same matters likewise written in Greek, to be carried to you. For if I had read yours first, you would say I had stolen from you. And though yours --- I read it gladly --- seemed to me a little rough and uncombed, yet it was for that very reason adorned, since it had neglected adornment, and like women who, by smelling of nothing, seem to smell sweetly. My book, however, has used up the whole perfumery-box of Isocrates and all his pupils' little caskets, and not a few of Aristotle's pigments too. You, at Corcyra, as another letter of yours signals, took only a quick look at it; afterwards, I take it, you got it from Cossinius. Which book I should not have dared to send you if I had not lazily and fastidiously approved it.

\ciceroSection{2} Although Posidonius has now written back from Rhodes that, when he read that memoir of ours which I had sent to him to write more elaborately on the same matters, he was not only not stirred to writing but plainly deterred. What more can I say? I have thrown the Greek nation into confusion. So those who were pressing me on every side to give them something to adorn have already stopped giving me trouble. You, if you like the book, will see to it that it be at Athens and in the other towns of Greece; for it seems likely to bring some light to my affairs.

\ciceroSection{3} As for the little speeches, both those you ask for and more besides, I shall send. For since the things which I write, stirred up by the zeal of young men, please you also: it suited me, since in those speeches which are called Philippics that fellow-citizen of yours Demosthenes had shown his brilliance, and had cut himself off from this stiffer judicial style of speaking that he might appear more dignified \textit{[Greek: semnoteros]} and more political \textit{[Greek: politik\=oteros]}, to take care that I should also have speeches that were called consular. Of which one is in the senate on the Kalends of January, the second to the people on the agrarian law, the third on Otho, the fourth for Rabirius, the fifth on the children of the proscribed, the sixth when I gave up my province in a public meeting, the seventh when I dismissed Catiline, the eighth which I delivered to the people on the day after Catiline fled, the ninth in the meeting on the day the Allobroges turned informer, the tenth in the senate on the Nones of December. There are besides two short pieces, like fragments \textit{[Greek: apospasmatia]} of the agrarian law. This whole body \textit{[Greek: s\=oma]} I will see that you have. And since both my writings and my deeds delight you, from those same books you will perceive both what I have done and what I have said. Or you should not have asked; for I was not pressing myself on you.

\ciceroSection{4} As for what you ask --- why I summon you, and at the same time signal that you are entangled in business and yet do not refuse, not only if there is need but even if I should wish, to come at once --- there is no real necessity, but yet you seemed to me able to lay out the times of your travels more conveniently. You are away too long, especially since you are in nearby places, and I do not enjoy you and you go without us. And just now there is peace; but if the madness of Pulchellus could go forward a little, I should rouse you from those parts vehemently. But Metellus admirably hinders him and will hinder him. What more? He is a consul who loves his country \textit{[Greek: philopatris]} and, as I have always judged, naturally good.

\ciceroSection{5} The other man, however, does not pretend, but plainly desires to be made tribune of the plebs. When the matter was being acted on in the senate, I broke the man and rebuked his inconsistency, who at Rome was seeking the tribunate of the plebs when he had been declaring in Sicily that he was seeking an inheritance --- and I said that we did not have to labour greatly, since it would no more be permitted him as a plebeian to ruin the commonwealth than it had been permitted to patricians like him while I was consul. Now when he had said that he had come from the Strait in seven days, that no one had been able to come out and meet him, and that he had entered by night, and was boasting of this in the assembly, I said that nothing new had happened to him: ``From Sicily to Rome in seven days. Three hours before that, from Rome to Interamna. You entered by night --- the same as before. No one came out to meet you --- not even then, when they most ought to have come.'' What more? I am making the impudent man modest, not only by the unbroken gravity of my speech but also by this sort of jokes. So now I rally and joke familiarly with him; indeed, when we were escorting a candidate, he asked me whether I had been wont to give the Sicilians a place at the gladiators. I said no. ``But I,'' he said, ``the new patron, will set the practice up. But my sister, who has so much consular space, gives me only one foot of room.'' --- ``Don't complain,'' I said, ``of one foot of your sister's; you are permitted to lift the other one too.'' Not a consular saying, you will say. I confess; but I detest that woman, the wretched consular's wife. For she is a seditious woman, who makes war with her husband --- and not only with Metellus but also with Fabius, because she takes ill that they ``are this in this'' [text corrupt].

\ciceroSection{6} What you ask about the agrarian law: it now seems to have cooled down. As to your gently rebuking me for my familiarity with Pompey: do not suppose that I have joined myself with him for my own protection, but the matter was so set up that, if there were by chance any disagreement between us, the greatest discords in the commonwealth would be inevitable. This has been forestalled and provided against by me, not so that I should give up that best line of mine, but so that he should be the better and lay aside something of his popular levity. Of my deeds, into which many had stirred him, know that he speaks more gloriously than of his own; for to himself he attributes the well-doing, to me the saving of the commonwealth. How much it profits me that he does this I do not know; certainly it profits the commonwealth. What if I were even to make Caesar better, whose winds are now greatly favourable --- am I doing so much harm to the commonwealth?

\ciceroSection{7} Indeed, even if no one envied me, if all favoured me, as was fair, yet that medicine which heals the diseased parts of the commonwealth must be no less approved than that which cuts them out. Now in fact, when that cavalry whom I, with you as standard-bearer and chief, had stationed on the Capitoline slope, has deserted the senate; and when our principal men think they touch the sky with a finger if there are bearded mullets in their fish-ponds that come to their hand, while they neglect the rest --- do I not seem to do enough good if I bring it about that those who can do harm do not wish to?

\ciceroSection{8} For our Cato --- you do not love him more than I do; but yet, with the best mind and the highest faith, he sometimes harms the commonwealth. He gives his opinion as if he were in the Republic of Plato \textit{[Greek: en t\=ei Plat\=onos politeiai]}, not in the dregs of Romulus. What is truer than that the man who took money for judging should come to trial? Cato moved this; the senate agreed; the equestrians made war on the senate-house, not on me (for I disagreed). What more impudent than the publicans renouncing their contract? Yet some sacrifice had to be made for the sake of keeping that order. Cato withstood and prevailed. So now, when a consul has been shut up in prison, when sedition has been stirred often, no one of those by whose flocking I, and likewise the consuls who came after me, were wont to defend the commonwealth, has come forward. ``What then?'' you will say; ``shall we have these men hired by wages?'' What shall we do, if otherwise we cannot? Or shall we be slaves to freedmen and to slaves? But, as you say, ``enough of zeal'' \textit{[Greek: halis spoud\=es]}.

\ciceroSection{9} Favonius carried my tribe more honourably than his own, lost Lucceius's. He prosecuted Nasica, dishonourably enough but with moderation; he spoke as though at Rhodes he had given his pains rather to the millers than to Molon. With me, because I had defended him, he was lightly displeased. Now nevertheless he is again a candidate for the commonwealth's sake. What Lucceius is doing I shall write to you when I have seen Caesar, who will be here in two days.

\ciceroSection{10} As to the Sicyonians wronging you, you attribute it to Cato and his rival Servilius. Well: does that blow not bear on many good men? But if so we please, let us praise them, and then in dissensions be left alone?

\ciceroSection{11} My Amaltheum looks for you and needs you. The Tusculan and Pompeian villas delight me greatly --- except that they have buried me, the very avenger of debt, in debt; not Corinthian bronze but this Forum brass. In Gaul we hope there is peace. Expect my Prognostica together with the little speeches in a few days; and yet, what you think about your arrival, write to me. For Pomponia bade word be brought to me that you would be at Rome in the month of Quintilis. That disagreed with the letter you had sent me about your census.

\ciceroSection{12} Paetus, as I wrote to you before, has given me all the books that his brother had left. This gift of his depends on your diligence. If you love me, see that they are kept and brought to me. Nothing can please me more than this. Both the Greek books and especially, with diligence, the Latin --- I should like you to keep. I shall reckon this as a little gift of yours. To Octavius I have given a letter; with him I had said nothing; for I did not think those provincial businesses of yours were a matter for him, nor did I count you among the petty money-lenders. But I wrote, as I should, with care.
